\documentclass{article}
\usepackage[utf8]{inputenc}
\usepackage{graphicx}
\usepackage{amsmath}
\usepackage{bm}

\title{2D Diffusion - Linear Case}
\author{The Weapon Group}
\date{December 2025}

\begin{document}

\maketitle

\section{Basic Setup - Dimensionless Variables and Discretization}
The 2 new equations:
\begin{equation}
\frac{\partial E(\rho, x,\tau)}{\partial \tau} - \frac{\partial^2E(\rho, x,\tau)}{\partial \rho^2} - \frac{1}{\rho}\frac{\partial E(\rho, x,\tau)}{\partial \rho} - \frac{\partial^2E( \rho, x,\tau)}{\partial x^2} = U(\rho, x,\tau) - E(\rho, x,\tau),
\end{equation}

\begin{equation}
\frac{\partial U(\rho, x,\tau)}{\partial \tau} = E(\rho, x,\tau) - U(\rho, x,\tau),
\end{equation}
Now lets discretize the equations:

\begin{equation}
\begin{split}
\frac{E^{n+1}_{i,j}-E^{n}_{i,j}}{\Delta \tau} & - \frac{E^{n+1}_{i,j+1}-2E^{n+1}_{i,j}+E^{n+1}_{i,j-1}}{\Delta \rho^2} 
- \frac{1}{\rho_j}\frac{E^{n+1}_{i,j+1}-E^{n+1}_{i,j-1}}{2\Delta \rho} \\
& - \frac{E^{n+1}_{i+1,j}-2E^{n+1}_{i,j}+E^{n+1}_{i-1,j}}{\Delta x^2} = U^{n+1}_{i,j} - E^{n+1}_{i,j},
\end{split}
\end{equation}

\begin{equation}
\frac{U^{n+1}_{i,j}-U^{n}_{i,j}}{\Delta \tau} = E^{n+1}_{i,j} - U^{n+1}_{i,j},
\end{equation}

Rearranging the equations to solve for the next time step values:
\begin{equation}
\begin{split}
E^{n+1}_{i,j} = E^{n}_{i,j} + \Delta \tau \Biggl( & \frac{E^{n+1}_{i,j+1}-2E^{n+1}_{i,j}+E^{n+1}_{i,j-1}}{\Delta \rho^2} 
+ \frac{1}{\rho_j}\frac{E^{n+1}_{i,j+1}-E^{n+1}_{i,j-1}}{2\Delta \rho} \\
& + \frac{E^{n+1}_{i+1,j}-2E^{n+1}_{i,j}+E^{n+1}_{i-1,j}}{\Delta x^2} + U^{n+1}_{i,j} - E^{n+1}_{i,j} \Biggr),
\end{split}
\end{equation}

\begin{equation}
U^{n+1}_{i,j} = \frac{U^{n}_{i,j} + \Delta \tau E^{n+1}_{i,j}}{1 + \Delta \tau},
\end{equation}

Define $\alpha$, $\lambda_x$, and $\lambda_\rho$ as:
\begin{equation}
\alpha = \frac{\Delta \tau}{1+\Delta \tau}, \quad \lambda_x = \frac{\Delta \tau}{\Delta x^2}, 
\quad \lambda_{\rho} = \frac{\Delta \tau}{\Delta \rho^2}, \quad \gamma_j = \frac{\Delta \tau}{2j\Delta \rho^2} = \frac{\lambda_{\rho}}{2j},
\end{equation}
We can rewrite the equations as (substituting $U^{n+1}_{j,i}$ into the first equation):
\begin{equation}
\begin{split}
E^{n+1}_{i,j}-E^{n}_{i,j} & = \lambda_{\rho}(E^{n+1}_{i,j+1}-2E^{n+1}_{i,j}+E^{n+1}_{i,j-1}) 
 + \frac{\lambda_{\rho}}{2j}(E^{n+1}_{i,j+1}-E^{n+1}_{i,j-1}) \\
& + \lambda_x({E^{n+1}_{i+1,j}-2E^{n+1}_{i,j}+E^{n+1}_{i-1,j}}) + \alpha U^{n}_{i,j} -\alpha E^{n+1}_{i,j},
\end{split}
\end{equation}
Now we will have to rearrange to have the $E^{n+1}$ terms on the left side:
\begin{equation}
\begin{split}
(1 + 2\lambda_{\rho} + 2\lambda_x + \alpha)E^{n+1}_{i,j} & - (\lambda_{\rho} + \frac{\lambda_{\rho}}{2j})E^{n+1}_{i,j+1} - (\lambda_{\rho} - \frac{\lambda_{\rho}}{2j})E^{n+1}_{i,j-1} \\
& - \lambda_x E^{n+1}_{i+1,j} - \lambda_x E^{n+1}_{i-1,j} = \alpha U^{n}_{i,j} + E^{n}_{i,j},
\end{split}
\end{equation}
Now we will to look for a way to solve this system of equations.

\subsection*{Matrix Formulation}
We would like to express this in matrix form $A\cdot E^{n+1} = B$, where $A$ is 
the coefficients of the $E^{n+1}$ terms, and $B$ is the right side of the equation.
But first we need to linearize the 2D indices $(i,j)$, where j is the $\rho$ index and i is the $x$ index, into a 1D index $k$:
\begin{equation}
k = i\cdot N_\rho + j,
\end{equation}
Where $N_x$ is the number of grid points in the x direction, that will to $E_{i+1,j} = E_{(i+1)\cdot N_\rho + j} = E_{k+N_\rho}$.
On the other hand, $E_{i,j+1} = E_{i\cdot N_\rho + j + 1} = E_{k+1}$. Now we can express the equation in terms of $k$:
\begin{equation}
\begin{split}
(1 + 2\lambda_{\rho} + 2\lambda_x + \alpha)E^{n+1}_{k} & - (\lambda_{\rho} + \frac{\lambda_{\rho}}{2j})E^{n+1}_{k+1} - (\lambda_{\rho} - \frac{\lambda_{\rho}}{2j})E^{n+1}_{k-1} \\
& - \lambda_x E^{n+1}_{k+N_\rho} - \lambda_x E^{n+1}_{k-N_\rho} = \alpha U^{n}_{k} + E^{n}_{k},
\end{split}
\end{equation}
This system is far from be efficient to solve directly, so we will look for other methods, it has a runtime of $O(N^6)$.

\subsection*{Silvester Equations Approach}
We can express the system as a Sylvester equation of the form:
\begin{equation}
A X^{n+1} + X^{n+1} B = Y^{n},
\end{equation}
Where $A$ and $B$ are matrices representing the discretization in the $\rho$ and $x$ directions respectively, and $C$ is the right-hand side matrix.
To construct matrices $A$ and $B$, we can use the coefficients from the discretized equations. The matrix $A$ will represent the axial diffusion terms, while matrix $B$ will represent the radial diffusion terms.
The right-hand side matrix $Y^{n}$ will be constructed from the known values of $E^{n}$ and $U^{n}$ at the current time step.
This approach allows us to leverage efficient numerical algorithms for solving Sylvester equations, which can significantly reduce
the computational complexity compared to direct methods, such as the method we've seen in the last subsection.

\subsection*{Details on Silvester Equations in Our Problem}
We will define the matrix:
\begin{equation}
    X^{n+1}\in {R}^{N_x\times N_{\rho}}\;\;\;\;[X^{n+1}]=E_{i,j}^n
\end{equation}
Where we think about a small cross of variables inside $X^n$ that looks like the following:
\begin{equation}
    \begin{bmatrix}
        * & E_{i-1,j} & * \\
        E_{i,j-1} & E_{i,j} & E_{i,j+1} \\
        * & E_{i+1,j} & *
    \end{bmatrix}
\end{equation}
Following the discussion of the implicit solution in the simpler 1D diffusion problem, as described in the last subsection - one can notice that the 5-\_\_\_\_ equation can be written in matrix form as
\begin{equation}
    AX^{n+1}+X^{n+1}B=Y^n
\end{equation}
    \begin{equation}
\begin{split}
(1 + 2\lambda_{\rho} + 2\lambda_x + \alpha)E^{n+1}_{i,j} & - (\lambda_{\rho} + \frac{\lambda_{\rho}}{2j})E^{n+1}_{i,j+1} - (\lambda_{\rho} - \frac{\lambda_{\rho}}{2j})E^{n+1}_{i,j-1} \\
& - \lambda_x E^{n+1}_{i+1,j} - \lambda_x E^{n+1}_{i-1,j} = \alpha U^{n}_{i,j} + E^{n}_{i,j},
\end{split}
\end{equation}
\begin{equation}
    a_{ij}E_{i-1,j}+b_{ij}E_{i,j}+c_{ij}E_{i+1,j}+d_{ij}E_{i,j+1}+e_{ij}E_{i,j+1}+f_{ij}E_{i,j}^*
\end{equation}
(*) Note that $E_{i,j}^*=2\lambda_\rho$  is the part of the center of the cross element and will appear in the $B$ matrix.
where the matrix $A,B,Y^n$ are defined as
\begin{gather}
A = 
\begin{bmatrix}
1 & 1 & 0 & \cdots & 0 & 0\\
a_1 & b_1 & c_1 & \cdots & 0 & 0 \\  
0 & a_2 & b_2 & c_2 &\cdots & 0 \\
\vdots & \vdots & \vdots & \ddots & \vdots & \vdots \\
0 & 0 & \cdots & a_{N_\rho-2} & b_{N_\rho-2} & c_{N_\rho-2}  \\
0 & 0 & 0 & 0 & \cdots & 1
\end{bmatrix}_{N_x\times N_x} = \\
\begin{bmatrix}
1 & 1 & 0 & \cdots & 0 & 0\\
-\lambda_x & 1+2\lambda_x+\alpha & -\lambda_x & \cdots & 0 & 0 \\  
0 & -\lambda_x &  1+2\lambda_x+\alpha & -\lambda_x &\cdots & 0 \\
\vdots & \vdots & \vdots & \ddots & \vdots & \vdots \\
0 & 0 & \cdots & -\lambda_x &  1+2\lambda_x+\alpha & -\lambda_x  \\
0 & 0 & 0 & 0 & \cdots & 1
\end{bmatrix}_{N_x\times N_x}
\end{gather}
and
\begin{gather}
B = 
\begin{bmatrix}
1 & d_1 & 0 & \cdots  & 0 & 0\\
1 & e_1  & d_2 & \cdots & 0 &0 \\  
0 & f_1 &  e_2 &\vdots & 0 \\
\vdots & \vdots & f_2 & \ddots & d_{N_x-2} & 0 \\
0 & 0 & \vdots & \cdots &  e_{N_x-2} & \vdots  \\
0 & 0 & 0 & \cdots &  f_{N_x-2} & 1
\end{bmatrix}_{N_{\rho}\times N_{\rho}}
\end{gather}
\begin{gather}
\begin{bmatrix}
1 & -\lambda_\rho(1-\frac{1}{2j}) & 0 & \cdots  & 0 & 0\\
1 & 2\lambda_\rho  & -\lambda_\rho(1-\frac{1}{2j}) & \cdots & 0 &0 \\  
0 & -\lambda_\rho(1+\frac{1}{2j}) &  2\lambda_\rho & \cdots &\vdots & 0 \\
\vdots & \vdots & -\lambda_\rho(1+\frac{1}{2j}) & \ddots & -\lambda_\rho(1-\frac{1}{2j}) & 0 \\
0 & 0 & \vdots & \cdots &  2\lambda_\rho & \vdots  \\
0 & 0 & 0 & \cdots & -\lambda_\rho(1+\frac{1}{2j}) & 1
\end{bmatrix}_{N_{\rho}\times N_{\rho}}
\end{gather}
The boundary conditions for this problem is
\ \begin{gather}
    E(x=0,\rho,t)=a(T_0^z)^4 \\
    E(x,\rho=0,t)=E(x,\rho=\Delta \rho,t) \\
\end{gather}
and the initial conditions are
\begin{equation}
    E(x,\rho,t=0)=U(x,\rho,t=0)=a(T_*)^4
\end{equation}

\section{The Dimension full Problem}
The dimensional problem is given by the following equations:
\begin{equation}
\frac{\partial E(r,z,t)}{\partial t} - \frac{c}{3\sigma}(\frac{\partial ^2 E}{\partial r^2} + 
\frac{1}{r}\frac{\partial E}{\partial r} + \frac{\partial ^2 E}{\partial z^2}) = c\sigma(U - E)
\end{equation}
\begin{equation}
\frac{1}{\varepsilon}\frac{\partial U(r,z,t)}{\partial t} = c\sigma(E - U)
\end{equation}

Now we will Derive the Numerical Scheme for the dimensional problem, similar to the dimensionless one.
\subsection*{Numerical Scheme for the Dimensional Problem}
First we discretize the equations:
\begin{equation}
\begin{split}
\frac{E^{n+1}_{i,j}-E^{n}_{i,j}}{\Delta t} & - \frac{c}{3\sigma} \Biggl( \frac{E^{n+1}_{i,j+1}-2E^{n+1}_{i,j}+E^{n+1}_{i,j-1}}{\Delta r^2}
    + \frac{1}{r_j}\frac{E^{n+1}_{i,j+1}-E^{n+1}_{i,j-1}}{2\Delta r} \\
& + \frac{E^{n+1}_{i+1,j}-2E^{n+1}_{i,j}+E^{n+1}_{i-1,j}}{\Delta z^2} \Biggr) = c\sigma(U^{n+1}_{i,j} - E^{n+1}_{i,j}),
\end{split}
\end{equation}
\begin{equation}
\frac{1}{\varepsilon}\frac{U^{n+1}_{i,j}-U^{n}_{i,j}}{\Delta t} = c\sigma(E^{n+1}_{i,j} - U^{n+1}_{i,j}),
\end{equation}
Rearranging the equations to solve for the next time step values:
\begin{equation}
\begin{split}
E^{n+1}_{i,j} = E^{n}_{i,j} + \Delta t \Biggl( & \frac{c}{3\sigma} \Biggl( \frac{E^{n+1}_{i,j+1}-2E^{n+1}_{i,j}+E^{n+1}_{i,j-1}}{\Delta r^2}
    + \frac{1}{r_j}\frac{E^{n+1}_{i,j+1}-E^{n+1}_{i,j-1}}{2\Delta r} \\
& + \frac{E^{n+1}_{i+1,j}-2E^{n+1}_{i,j}+E^{n+1}_{i-1,j}}{\Delta z^2} \Biggr) + c\sigma(U^{n+1}_{i,j} - E^{n+1}_{i,j}) \Biggr),
\end{split}
\end{equation}
\begin{equation}
U^{n+1}_{i,j} = \frac{U^{n}_{i,j} + c\sigma \varepsilon \Delta t E^{n+1}_{i,j}}{1 + c\sigma \varepsilon \Delta t},
\end{equation}
Define the following parameters:
\begin{equation}
f = \frac{1}{1+c\sigma \varepsilon \Delta t}, \lambda_r = \frac{c\Delta t}{3\sigma \Delta r^2}, \lambda_z = \frac{c\Delta t}{3\sigma \Delta z^2}, 
\gamma_j = \frac{c\Delta t}{3\sigma 2r_j \Delta r},
\end{equation}
Notice that $r_j = j\Delta r$, one can substitute it back to the equations.
\begin{equation}
    \gamma_j = \frac{\lambda_r}{2j}
\end{equation}
We can rewrite the equations as (substituting $U^{n+1}_{j,i}$ into the first equation):
\begin{equation}
\begin{split}
E^{n+1}_{i,j}-E^{n}_{i,j} & = \lambda_{r}(E^{n+1}_{i,j+1}-2E^{n+1}_{i,j}+E^{n+1}_{i,j-1}) 
 + \frac{\lambda_r}{2j}(E^{n+1}_{i,j+1}-E^{n+1}_{i,j-1}) \\
& + \lambda_z({E^{n+1}_{i+1,j}-2E^{n+1}_{i,j}+E^{n+1}_{i-1,j}}) + c\sigma (1-f) U^{n}_{i,j} - c \sigma (1-f) E^{n+1}_{i,j},
\end{split}
\end{equation}
Rearranging to have the $E^{n+1}$ terms on the left side:
\begin{equation}    
\begin{split}
(1 + 2\lambda_{r} + 2\lambda_z + c\sigma (1-f))E^{n+1}_{i,j} & - (\lambda_{r} + \frac{\lambda_{r}}{2j})E^{n+1}_{i,j+1} 
- (\lambda_{r} - \frac{\lambda_{r}}{2j})E^{n+1}_{i,j-1} \\
& - \lambda_z E^{n+1}_{i+1,j} - \lambda_z E^{n+1}_{i-1,j} = c\sigma (1-f) U^{n}_{i,j} + E^{n}_{i,j},
\end{split}
\end{equation}
This system can be solved using the Sylvester equation approach as described in the dimensionless case.



\end{document}